\documentclass[a4paper,14pt]{exam}
\usepackage{array,amsmath,fancybox,amsfonts,bm}
\usepackage{graphicx,amssymb,wasysym,latexsym,slashed,anysize,subfigure,calc}



%versores: %uso: \uveci\ne\hat{\imath}_{\uvec{i}}+\uvec{j}+\uvec{k}$
	\usepackage{bm}
	\newcommand{\uveci}{{\bm{\hat{\textnormal{\bfseries\i}}}}}
	\newcommand{\uvecj}{{\bm{\hat{\textnormal{\bfseries\j}}}}}
	\DeclareRobustCommand{\uvec}[1]{{%
 	 \ifcsname uvec#1\endcsname
  	   \csname uvec#1\endcsname
  	 \else
   		 \bm{\hat{\mathbf{#1}}}%
   		\fi
									}}



%tamaños fuentes extra: 8,9,14,17,20
	\usepackage{extsizes}
	
 %caja bonita
  \usepackage[many]{tcolorbox}
\definecolor{titlebg}{RGB}{80,100,72}
\newtcolorbox{learn}{
  breakable,
  enhanced,
  arc=0pt,
  outer arc=0pt,
  colframe=titlebg,
  colback=titlebg!7,
  overlay unbroken and first={
    \node[
      draw=titlebg,
      fill=titlebg,
      rotate=270,
      anchor=north west,
      text=white,
      font=\bfseries
    ]
    at (frame.north west)  
    {DATOS};
  }
}

%imágenes fundidas texto
\usepackage{wrapfig}

\usepackage{mwe,float,afterpage}

\usepackage{anysize,float}
	%\marginsize{left}{right}{top}{bottom}
    \marginsize{1.5cm}{1.5cm}{1cm}{.7cm}
    
    
    
    
%entorno para química: texto y math mode
%\newcommand*\chem[1]{\textup{#1}}
\newcommand*\chem[1]{\ensuremath{\mathrm{#1}}}

%símbolo bonito para qed
\usepackage{amsthm}
	%\newcommand{\mc}[3]{\multicolumn{#1}{#2}{#3}}
\newcommand{\halmos}{\vspace*{-7mm}
	
					\begin{flushright}
						\rule{1mm}{2.5mm} 
					\end{flushright}\vspace*{-6mm}
						
					  }
\renewcommand{\qed}{\halmos}

\usepackage[spanish]{babel}
\usepackage[utf8]{inputenc}
%\usepackage{fancyhdr}
\usepackage{lettrine}

%subrayado y highlight. Tiene que estar detrás de usepackage[color]
\usepackage{soul}
	%\textst{normal} 
	\DeclareRobustCommand{\hlcyan}[1]{{\sethlcolor{cyan}\hl{#1}}} %en cyan!
		%\hl{highlighted text} %uso


%fuentes
	\usepackage[T1]{fontenc}
	\usepackage{mathpazo}

	\usepackage{librebaskerville}

%fancy tables
	\usepackage{booktabs}
	\usepackage[small,centerlast,up]{caption}

	

\usepackage{amssymb,amsmath}
\newcommand{\angstrom}{\mbox{\normalfont\AA}}

\usepackage{graphicx,graphics}
\usepackage{color}
	%	\colorgrids %grids con color
	%	\definecolor{GridColor}{gray}{0.8}
		%grid
%			\setlength{\gridsize}{5mm}
%   			 \setlength{\gridlinewidth}{0.1pt}

%fracciones en texto más elaboradas
	\usepackage{xfrac}

%mejora en la exposición de guarismos
	\usepackage{numprint}
 


%para los entornos de soluciones
\usepackage{mdframed,xcolor,color}
	%uso: 	entorno (begin,end): {mdframed}[backgroundcolor=brown!10] 


\usepackage{everypage,draftwatermark} %Marca de agua
\SetWatermarkText{\parbox{2\textwidth}{\centering B2\\Dist}}
\SetWatermarkScale{1.5}
\SetWatermarkLightness{0.92}





\pagestyle{headandfoot}
%\runningheadrule
%\firstpageheader{Física para Biólogos}{Boletín 1a}{1 de octubre de 2011}
%\runningheader{Física para Biólogos}
%              {Boletín 1a, Página \thepage\ de \numpages}
%            {1 de octubre de 2011}
%\firstpagefooter{}{}{}
%\runningfooter{}{}{}


%espacios extra
%\extrawidth{1in}
%%\extrafootheight{.75in}
%\extraheadheight{1in}

%cabecera
\header{{\small DOSSIER DE PROBLEMAS}
		}       
      {%\\Noviembre de 2019
       }
       {B2-distancia}
       
%pie de página

   %enumera en todas menos en la primera
   %\runningfootrule

   %enumera en todas
   %\footrule
   %\footer{}{Página \thepage\ de \numpages}{}
%   \footer{}
%	{Página \thepage\ de \numpages}
%	{\small\bfseries\iflastpage{\textsc{Fin del examen}}{%Da la vuelta a la página\ldots
%	}}

       

		
\begin{document}
	


\begin{enumerate}
	
		%ondas banda ap
		\item La perturbación asociada a una onda viene descrita por la expresión
$$\psi(z,t)= 10^{-3}\cos\left(12t +3\pi z-\frac{3\pi}{4}\right),\; \text{\small (SI)}.$$
	\begin{itemize}
		\item[a)] Determina: \textit{el sentido y la dirección de propagación, $\lambda$ y $c$}.
		\item[b)] Obtén la \textit{diferencia de fase}, en el mismo instante de tiempo, para dos puntos separados una distancia igual a $3\lambda$.
	\end{itemize}
	\qed
		
		%sonido banda ap.
		\item 
			\begin{itemize}
				\item[a)] El sonido emitido por un altavoz es registrado con un nivel sonoro de $70\;dB$ a una 
distancia de $2\;m$. Si el altavoz se aproxima como un foco puntual, estima la \textit{distancia} a la que el sonido deja de ser perceptible para un ser humano.
				\item[b)] Una persona hablando produce un sonido con una intensidad unas quinientas veces mayor que 						cuando susurra. ¿Qué \textit{diferencia de nivel sonoro} hay entre ambos sonidos?
				\vspace*{.2mm}					
			\end{itemize}
		\qed
		
		
		%banda notable
		\item 	
			\begin{itemize}
				\item[a)] Una lámpara de luz monocromática de frecuencia $f=5\cdot 10^{15}\,Hz$ emite con $20\;W$ de potencia.  Suponiendo que la radiación se propaga en el vacío, determina su \textit{longitud de onda} y el \textit{número de fotones que se emiten por segundo}.
				\item[b)] Las salanganas %y \textit{vencejos de cueva} 
			son pájaros tropicales que vuelan en completa oscuridad gracias al fenómeno de ecolocalización: pueden detectar la difracción producida por obstáculos tras emitir ondas sonoras sobre ellos. Uno de estos pájaros puede emitir con una frecuencia máxima de $8\cdot 10^3\; Hz$, ¿cuál es el \textit{orden de magnitud de la longitud} del objeto más pequeño que puede detectar en su vuelo?
			\end{itemize}
		\qed
		
		%ondas banda not-sob.
		\item En un estadio de fútbol es común que el público haga «la ola» si está de buen rollo. Los espectadores levantan los brazos cada $10\;s$ y la ola producida es tal que cualesquiera dos de ellos sentados en la misma fila y separados por una distancia de $50\;m$ se mueven sincrónicamente. Modelizando esta propagación  como una onda armónica:
		\begin{itemize}
		\item[a)] ¿Qué \textit{tipo de polarización} presenta? Calcula su \textit{longitud de onda}.			
		\item[b)] Las manos de cada espectador, al pasar la ola y levantarse, se mueven verticalmente $1\;m$ desde el punto más bajo al más alto. Escribe una posible \textit{función} para la onda que se propaga en una misma fila asumiendo que $y(0,0)=0$.
		\end{itemize}
		\qed
		
\end{enumerate}		
 
\end{document}

